% SIGMOD 2027 研究论文投稿。格式依据官方 CFP：
%   https://2027.sigmod.org/calls_papers_sigmod_research.shtml
%   - 正文 ≤ 12 页，参考文献不计入且页数不限（录用后 camera-ready 为 13 页且不得有附录）
%   - 双盲：正文任何位置不得出现作者、单位、致谢、资助来源
%   - 2 栏 ACM 格式，字号/页边距/栏间距/行距一律不得改动
%   - 提交 PDF ≤ 10 MB，Letter 纸
%
% 类选项：sigconf = ACM 会议 2 栏格式；review = 给正文加行号方便审稿人引用；
% anonymous = 隐去作者与单位。定稿时把 review,anonymous 两项删掉即可。
\documentclass[sigconf,review,anonymous]{acmart}

% —— 投稿期去掉版权块 ——
% acmart 默认会在首页脚印 ACM 版权声明和 DOI，那些数字要等录用后从 ACM eRights
% 系统拿。投稿阶段用下面三行清掉；这只影响首页脚注区，不改动任何版面度量，
% 不违反 CFP 的"格式不得改动"。录用后连同 \acmConference 一起换成 eRights 给的内容。
\setcopyright{none}
\settopmatter{printacmref=false}
\renewcommand\footnotetextcopyrightpermission[1]{}

% —— camera-ready 时启用，内容以 ACM eRights 邮件为准 ——
% \copyrightyear{2027}
% \acmYear{2027}
% \setcopyright{acmlicensed}
% \acmConference[SIGMOD '27]{International Conference on Management of Data}{May 31--June 4, 2027}{Huntington Beach, CA, USA}
% \acmDOI{10.1145/nnnnnnn.nnnnnnn}
% \acmISBN{978-x-xxxx-xxxx-x/YY/MM}

\begin{document}

\title{Title Goes Here}

% anonymous 选项会自动把下面这些替换成 "Anonymous Author(s)"，所以真名可以先写着，
% 不必等到定稿再填。但正文、致谢、图表、URL、自引措辞里的身份线索必须自己管。
\author{Author One}
\affiliation{%
  \institution{Institution}
  \city{City}
  \country{Country}}
\email{one@example.com}

\renewcommand{\shortauthors}{Anonymous et al.}

\begin{abstract}
Abstract goes here.
\end{abstract}

% CCS 概念与关键词：投稿时非必需，camera-ready 必填。
% CCS 代码到 https://dl.acm.org/ccs 用 CCS 工具生成后粘贴回来。
% \begin{CCSXML} ... \end{CCSXML}
% \ccsdesc[500]{Information systems~Information extraction}
% \keywords{keyword one, keyword two}

\maketitle

\section{Introduction}
\label{sec:intro}

% Synchronized with introduction.zh.md on 2026-09-13.
% Narrative: agent activities -> properties of scholarly materials and knowledge
% needs -> research question -> CS empirical setting -> scenario and challenges
% -> P4A artifact organization and construction -> evaluation and significance.
% Papers and associated materials are the research objects. Specific organizing
% units and access paths belong in the methods section.
% A1--A4 denote activity types (Activity), which can be interleaved.
% Working draft: the query, candidate papers, and decision changes still require
% verification. Methods and evaluation remain under discussion; the research
% agenda below is not a list of completed contributions.
Agents assisting researchers with discovering related work, comparing research
approaches, and conducting experiments need to draw on knowledge from papers and
their associated materials. A study's problem, methods, experimental settings,
and results are distributed across its main text, figures, tables, and
appendices; supporting materials may provide further implementation and
evaluation details. Studies are also connected through citations, shared
resources, and the reuse of methods. To make justified decisions, agents need
to understand both the individual pieces of information and how they relate
within a particular study.

These materials are primarily organized to present and substantiate research.
The relationships among claims, supporting evidence, and applicability
conditions are not always explicitly represented as independently accessible
knowledge records. For example, determining whether a method applies to a
current problem may require relating its description in the main text to
results in tables, settings in an appendix, and an implementation in supporting
materials. Even after locating relevant content, an agent must establish
whether these pieces concern the same experiment, which conditions a conclusion
depends on, and whether the evidence supports the current choice. Interpreting
and checking these relationships across content and materials is part of the
work needed to organize existing research into knowledge that supports
subsequent research actions.

We therefore ask: \textbf{How can papers and their associated materials be
transformed into knowledge artifacts that agents can access, verify, and reuse
as needed?} Such artifacts need to organize problem definitions, methods,
experimental conditions, research conclusions, and associated resources into
records with explicit semantics, while preserving their relationships, source
evidence, and applicability conditions. Agents can use these records to
discover and screen candidates, trace and cite supporting evidence, and obtain
the knowledge needed to reproduce and use existing methods or resources. They
can return to the original sources when further inspection is needed.
We refer to these research
artifacts prepared for agent use as \emph{agent-ready artifacts}.

Computer science provides a suitable empirical setting for studying this
transformation. Preprints and open publication channels make recent research
accessible. Explicit task definitions, method descriptions, and experimental
settings provide cues for interpreting content, while public supporting
materials provide further evidence for inspecting implementations and
evaluations. In studies with these properties, code repositories, datasets,
and benchmarks supply verifiable information about implementation details,
data construction, and evaluation protocols that can be related to the paper's
account. We therefore focus our initial empirical study on computer science
papers with public supporting materials, using their accessibility and the
traceable relationships among research objects to study the construction and
verification of knowledge for agents. Even when these materials are public,
establishing whether they concern the same experiment, which conditions a
conclusion depends on, and whether a documented use supports a new need still
requires interpretation and assessment across materials.

\begin{quote}
\small
\textbf{Scenario: \texttt{<query>}.} An agent receives this request from a
researcher and needs to investigate related work and draw on existing research
to test an idea. We characterize its use of papers and associated materials
through four activity types, labeled A1--A4, where A stands for \emph{Activity}:

\textbf{A1---Discover.} Retrieve papers based on the researcher's needs to
form a candidate set of related work.\\[2pt]
\textbf{A2---Screen.} Refine the candidates using their problem definitions,
methods, and experimental conditions to decide whether to include them and
how to prioritize them.\\[2pt]
\textbf{A3---Expand and substantiate.} Follow citations, method reuse, and
resource relationships to find related work and materials, enrich the research
context, and cite studies and evidence supporting the current assessment or
research plan.\\[2pt]
\textbf{A4---Reproduce and use.} Use existing methods or resources in experiments,
including running implementations, reproducing results, or applying them to new
research tasks, while checking the required versions, configurations, and
evaluation conditions.
\end{quote}

% Scenario sketch, pending verification against actual materials:
% Initial retrieval returns <candidate paper A>. Inspecting <key evidence in the
% paper and associated materials> may reveal that its method or evaluation
% conditions differ from <the current need>, prompting a revised assessment.
% Further inspection may introduce <candidate paper B> through resource or
% citation relationships. Actual materials must determine the candidates,
% evidence, and decision changes. The final scenario should show which knowledge
% affected a choice and how it was obtained from the papers and materials.

As the scenario illustrates, an agent can assist a researcher through one or
more of A1--A4 in response to \texttt{<query>}. These activities can be
interleaved and affect one another: task examples or scoring code found through
A3 may change the screening outcome in A2, version differences or experimental
results observed in A4 may prompt reconsideration of an earlier judgment, and
new work found through citation and resource relationships may trigger further
retrieval (A1). Agents therefore need to
connect paper content with evidence from associated materials and interpret
their implications for subsequent research activities.

When other research needs involve the same paper or resource, its problem
definition, documented uses, and experimental conditions can be reused, while
applicability to a new need must be reassessed using these facts. Obtaining this
knowledge anew from the original materials incurs repeated retrieval, reading,
and reasoning costs. This creates a further need for knowledge organization:
preserving the results of interpretation and verification as accessible records,
together with their evidence and conditions, so that subsequent research
activities can reuse existing knowledge and continue inspection when needed.

Constructing such research artifacts presents three challenges. First,
\emph{establishing semantic relationships supported by evidence}. A paper's
methods, experimental conditions, and conclusions may appear in different
places, while implementation and evaluation materials supply further levels of
detail. To support screening in A2 and substantiation in A3, the system needs
to establish which experiments and results support a claim, what roles the
associated resources played, and which settings were used, avoiding
unsupported relationships assembled from individually correct pieces of
information. An author's general description, a documented use in a specific
experiment, and applicability to a current need must be represented separately,
with uncertainty retained where the evidence is insufficient.

Second, \emph{organizing knowledge units and relationships to support different
research activities}. A1--A4 require information at different levels of
granularity: discovery may begin with a problem or method, screening requires
conditions and evidence, expansion and substantiation require research
relationships and supporting sources, and reproduction and use call for
specific versions, configurations, and evaluation protocols.
A knowledge representation for agents needs to make clear
which content can be accessed independently, which relationships must be
preserved together, and how a record can be traced to the paper and associated
materials. Integration across papers must also distinguish different versions
and uses of resources with the same name, so that existing knowledge can
support subsequent retrieval and comparison.

Third, \emph{preserving the conditions for verifiable reuse}. Whether a record
provides sufficient support for a current action depends on its source,
version, scope of inspection, and the research need. A paper's claim that code
has been released, observed repository contents, successful program execution,
and experimental reproduction are distinct findings. Their evidential status
needs to be represented so that an agent can decide when to use an existing
record and when further checks are needed, without treating a conclusion
established under specific conditions as a general fact.

We study \emph{Papers for Agents (P4A)}, exploring how to organize and construct
research artifacts for agents. P4A aims to integrate research content from
papers and their publicly available associated materials into grounded,
accessible, and reusable knowledge records through preprocessing that includes
content parsing, semantic interpretation, and verification against source
materials, retaining these artifacts in a persistent literature resource
library. We investigate reasoning-based
extraction for establishing semantic relationships within papers and across
associated materials, organizing facts together with source evidence,
verification status, and applicability conditions. This would allow agents to
access existing knowledge as needed and pursue unresolved questions. The
granularity of records, their organization, and access mechanisms remain to be
determined for specific tasks and evidence environments.

We plan to pursue this work in three directions:
\begin{itemize}
\item \textbf{Organizing research artifacts for agents.} Investigate how to
represent research content, associated materials, and their relationships while
preserving sources, conditions, and verification status to support the four
activity types above.
\item \textbf{Constructing artifacts through reasoning and verification.}
Investigate how to extract, align, and verify semantic relationships across
papers and associated materials, handling insufficient or inconsistent evidence
to produce records that support subsequent research actions.
\item \textbf{Evaluating artifact quality and usefulness.} Independently assess
semantic correctness, relationship correctness, and evidential support, then
examine whether agents using these artifacts make better-supported research
choices and whether knowledge reuse reduces repeated reading and interpretation.
\end{itemize}

We plan to evaluate P4A on defined research tasks using a collection of computer
science papers and associated materials with evidence that can be independently
checked. The running scenario will illustrate how a knowledge artifact
affects a specific judgment, while the full set of tasks will examine whether
its benefits extend to other research needs. Method comparisons should provide
equivalent access to the original materials and declare the available budgets.
Baselines such as direct reading and retrieval and conventional structured
extraction will help distinguish the effects of organization, reasoning, and
verification. Comparisons with external DeepResearch systems under open search
can provide supplementary evidence, with each system's information access
documented. Construction and usage costs will be reported separately, and reuse
will be examined across multiple needs involving shared papers and resources.
The task scope, comparison protocols, and quality criteria remain to be defined.

This work explores ways of organizing published research to make it easier
for agents to discover, understand, and use. By transforming content distributed
across papers and associated materials into accessible, traceable research
knowledge with explicit applicability conditions, we aim to provide a reusable
foundation for agents assisting researchers with subsequent work. Its
effectiveness must be evaluated through both artifact quality and the research
activities it supports.

% TODO: Select a representative query from the E07 paper pool, replace scenario
% placeholders with verified candidates and evidence, and retain the actual
% retrieval and assessment records.
% TODO: Support the CS empirical setting with literature or sample evidence;
% avoid generalizing properties of the selected public studies to the field.
% Establish specific gaps through related work and baseline observations.
% TODO: Define the initial task slice, knowledge units, and evidence boundaries.
% Specify construction and access mechanisms in later sections.
% TODO: Specify which activity types are evaluated and the extent of execution
% and reproduction. Define independent verification evidence, ground truth,
% quality criteria, and budgets. Separate exploratory scenarios from evaluation
% tasks and replace the research agenda with established contributions and
% measured findings, keeping the Chinese and English drafts aligned.

\section{Conclusion}
\label{sec:conclusion}

% 页数闸门的锚点：Makefile 的 `make check` 从 .aux 里读这个 label 所在页码，
% 判断正文是否超过 12 页。它必须紧挨在参考文献之前。
\label{endofbody}

% 双盲期不得出现致谢。录用后再取消注释。
% \begin{acks}
% Acknowledgements go here.
% \end{acks}

\bibliographystyle{ACM-Reference-Format}
\bibliography{refs}

\end{document}
